% Einheit 5 von 5 · Faktorisieren (binomische-formeln.md, Einheit 3) – Hauptnummern 39–46
\setcounter{aufgabe}{38}
\zweigkopf{Einheit 5 von 5 · Faktorisieren}{Hier lernst du, eine Summe mit einer binomischen Formel in ein Produkt zu verwandeln · neu in diesem Jahr (an der Oberschule je nach Buch bis Klasse 10) · keine P10-Aufgabe · baut auf: binomische Formeln (Einheit 4), Ausklammern (Einheit 2)}{e5}

\begin{aufgabe}{Ich erkenne, ob ein Term ein Quadrat ist. Trage ein, ob der Term das Quadrat einer Zahl oder eines Terms ist (ja oder nein). Wenn ja, schreibe dazu, wovon er das Quadrat ist. Faktorisiere nichts.}
\sachtabelle{lcc}{Term & Quadrat? & Quadrat von}{$81$ & \leerzelle & \leerzelle \\ $z^2$ & \leerzelle & \leerzelle \\ $25m^2$ & \leerzelle & \leerzelle \\ $8x$ & \leerzelle & \leerzelle \\ $12$ & \leerzelle & \leerzelle \\ $100y^2$ & \leerzelle & \leerzelle \\ $2x^2$ & \leerzelle & \leerzelle}
\end{aufgabe}

\begin{aufgabe}{Ich erkenne, ob das Mittelglied passt. Das erste und das letzte Glied sind Quadrate. Passt das Mittelglied zu einer binomischen Formel? Kreuze an. Faktorisiere nichts.}
\begin{teilezwei}
\tz $x^2 + 4x + 4$ \quad \janein & \tz $x^2 + 8x + 64$ \quad \janein \\
\tz $y^2 - 20y + 100$ \quad \janein & \tz $m^2 + 9m + 81$ \quad \janein \\
\tz $9x^2 + 6x + 1$ \quad \janein & \tz $z^2 + 10z + 36$ \quad \janein \\
\end{teilezwei}
\end{aufgabe}

\begin{aufgabe}{Ich kann eine Summe mit einer binomischen Formel in ein Produkt verwandeln. Schreibe als Produkt.}
\beispiel{x^2 - 36 &= (x + 6) \cdot (x - 6)}
\begin{gleichungsraster}[1]
\gl{x^2 - 16} & \gl{x^2 - 81} \\ \noalign{\vspace{5mm}}
\gl{m^2 - 1} & \gl{y^2 - 100} \\ \noalign{\vspace{5mm}}
\gl{x^2 + 16x + 64} & \gl{z^2 - 18z + 81} \\ \noalign{\vspace{5mm}}
\end{gleichungsraster}
\end{aufgabe}

\begin{aufgabe}{Ich kann eine Summe mit einer binomischen Formel in ein Produkt verwandeln – weiter. Schreibe als Produkt.}
\begin{gleichungsraster}[1]
\gl{16y^2 + 40y + 25} & \\ \noalign{\vspace{5mm}}
\end{gleichungsraster}
Klammere zuerst den gemeinsamen Faktor aus, dann schreibe als Produkt.
\begin{gleichungsraster}[2]
\gl{7x^2 - 28} & \\ \noalign{\vspace{5mm}}
\end{gleichungsraster}
\begin{teile}
\teil Prüfe, ob $x^2 + 6x + 16$ zu einer binomischen Formel passt. Wenn nicht, begründe.
\end{teile}
Schreibe als Produkt und mache die Probe durch Ausmultiplizieren.
\begin{gleichungsraster}[2]
\gl{16x^2 - 1} & \\ \noalign{\vspace{5mm}}
\end{gleichungsraster}
\end{aufgabe}

\begin{aufgabe}{Ich kann eine Summe mit einer binomischen Formel in ein Produkt verwandeln – weiter. Ein Produkt ist null, wenn einer seiner Faktoren null ist. Schreibe die linke Seite als Produkt und bestimme alle Lösungen.}
\begin{gleichungsraster}[3]
\gl{x^2 - 144 = 0} & \\ \noalign{\vspace{5mm}}
\end{gleichungsraster}
\begin{teile}
\teil Schreibe den Term als Quadrat einer Klammer plus eine Zahl.\\
$x^2 + 8x + 20 = (x + 4)^2 +$ \leerfeld
\end{teile}
Schreibe den Term vollständig als Produkt und mache die Probe durch Ausmultiplizieren.
\begin{gleichungsraster}[3]
\gl{2x^2 - 24x + 72} & \\ \noalign{\vspace{5mm}}
\end{gleichungsraster}
\end{aufgabe}

\begin{aufgabe}{Ich finde den Fehler beim Faktorisieren. Finde den Fehler, benenne ihn und korrigiere. Löse danach die Aufgabe darunter selbst: Schreibe als Produkt, wenn es geht.}
\begin{teile}
\teil Mia rechnet so:
\end{teile}
\rechnung{x^2 + 16 &= (x + 4)^2}
\begin{gleichungsraster}[1]
\gl{m^2 + 64} & \\ \noalign{\vspace{5mm}}
\end{gleichungsraster}
\begin{teile}
\teil Jan rechnet so:
\end{teile}
\rechnung{x^2 + 10x + 16 &= (x + 4)^2}
\begin{gleichungsraster}[1]
\gl{x^2 + 20x + 100} & \\ \noalign{\vspace{5mm}}
\end{gleichungsraster}
\begin{teile}
\teil Lina rechnet so:
\end{teile}
\rechnung{x^2 - 8x + 16 &= (x + 4)^2}
\begin{gleichungsraster}[1]
\gl{y^2 - 22y + 121} & \\ \noalign{\vspace{5mm}}
\end{gleichungsraster}
\end{aufgabe}

\begin{aufgabe}{Ich kann begründen, wann man mit einer binomischen Formel faktorisieren kann.}
\begin{teile}
\teil Begründe, warum man $z^2 + 36$ nicht mit einer binomischen Formel als Produkt schreiben kann.
\teil Kann man $9 - y^2$ als Produkt schreiben? Begründe.
\end{teile}
\end{aufgabe}

\begin{aufgabe}{Ich kann die Seitenlänge eines Quadrats aus seiner Fläche als Term bestimmen. Ein quadratisches Grundstück hat den Flächeninhalt $(w^2 + 24w + 144)$~m².}
\begin{teile}
\teil Gib die Seitenlänge des Grundstücks als Term an.
\teil Berechne die Seitenlänge und den Flächeninhalt für $w = 8$.
\teil Ein rechteckiges Nachbargrundstück hat den Flächeninhalt $(w^2 - 16)$~m². Eine Seite ist $(w + 4)$~m lang. Wie lang ist die andere Seite?
\end{teile}
\end{aufgabe}
